\documentclass[10pt,a4paper]{amsart}
\usepackage{amsmath,amssymb,amsthm,graphicx}
\usepackage{verbatim}
\newenvironment{reference}{\comment}{\endcomment}
%\newenvironment{reference}{}{}
\newenvironment{slogan}{\comment}{\endcomment}
\newenvironment{history}{\comment}{\endcomment}

% For commutative diagrams we use Xy-pic
\usepackage[all]{xy}

% We use 2cell for 2-commutative diagrams.
\xyoption{2cell}
% Theorem environments.
%
\theoremstyle{plain}
\newtheorem{theorem}[subsection]{Theorem}
\newtheorem{proposition}[subsection]{Proposition}
\newtheorem{lemma}[subsection]{Lemma}

\theoremstyle{definition}
\newtheorem{definition}[subsection]{Definition}
\newtheorem{example}[subsection]{Example}
\newtheorem{exercise}[subsection]{Exercise}
\newtheorem{situation}[subsection]{Situation}

\theoremstyle{remark}
\newtheorem{remark}[subsection]{Remark}
\newtheorem{remarks}[subsection]{Remarks}

\numberwithin{equation}{subsection}

% Macros
%
\def\lim{\mathop{\mathrm{lim}}\nolimits}
\def\colim{\mathop{\mathrm{colim}}\nolimits}
\def\Spec{\mathop{\mathrm{Spec}}}
\def\Hom{\mathop{\mathrm{Hom}}\nolimits}
\def\Ext{\mathop{\mathrm{Ext}}\nolimits}
\def\SheafHom{\mathop{\mathcal{H}\!\mathit{om}}\nolimits}
\def\SheafExt{\mathop{\mathcal{E}\!\mathit{xt}}\nolimits}
\def\Sch{\mathit{Sch}}
\def\Mor{\mathop{\mathrm{Mor}}\nolimits}
\def\Ob{\mathop{\mathrm{Ob}}\nolimits}
\def\Sh{\mathop{\mathit{Sh}}\nolimits}
\def\NL{\mathop{N\!L}\nolimits}
\def\CH{\mathop{\mathrm{CH}}\nolimits}
\def\proetale{{pro\text{-}\acute{e}tale}}
\def\etale{{\acute{e}tale}}
\def\QCoh{\mathit{QCoh}}
\def\Ker{\mathop{\mathrm{Ker}}}
\def\Im{\mathop{\mathrm{Im}}}
\def\Coker{\mathop{\mathrm{Coker}}}
\def\Coim{\mathop{\mathrm{Coim}}}

% Boxtimes
%
\DeclareMathSymbol{\boxtimes}{\mathbin}{AMSa}{"02}

%
% Macros for moduli stacks/spaces
%
\def\QCohstack{\mathcal{QC}\!\mathit{oh}}
\def\Cohstack{\mathcal{C}\!\mathit{oh}}
\def\Spacesstack{\mathcal{S}\!\mathit{paces}}
\def\Quotfunctor{\mathrm{Quot}}
\def\Hilbfunctor{\mathrm{Hilb}}
\def\Curvesstack{\mathcal{C}\!\mathit{urves}}
\def\Polarizedstack{\mathcal{P}\!\mathit{olarized}}
\def\Complexesstack{\mathcal{C}\!\mathit{omplexes}}
% \Pic is the operator that assigns to X its picard group, usage \Pic(X)
% \Picardstack_{X/B} denotes the Picard stack of X over B
% \Picardfunctor_{X/B} denotes the Picard functor of X over B
\def\Pic{\mathop{\mathrm{Pic}}\nolimits}
\def\Picardstack{\mathcal{P}\!\mathit{ic}}
\def\Picardfunctor{\mathrm{Pic}}
\def\Deformationcategory{\mathcal{D}\!\mathit{ef}}
\usepackage{cleveref}
\crefname{originaldefinition}{original definition}{original definitions}
\Crefname{originaldefinition}{Original definition}{Original definitions}
\makeatletter
\newlabel{definition-valuative-criterion}{{03IX}{source}{}{sourcepointer-1}{}}
\newlabel{definition-valuative-criterion@cref}{{[originaldefinition][0][]03IX}{[1][0][]source}{}{}{}}
\makeatother
\crefname{originallemma}{original lemma}{original lemmas}
\Crefname{originallemma}{Original lemma}{Original lemmas}
\makeatletter
\newlabel{lemma-quasi-compact-existence-universally-closed}{{03KA}{source}{}{sourcepointer-2}{}}
\newlabel{lemma-quasi-compact-existence-universally-closed@cref}{{[originallemma][0][]03KA}{[1][0][]source}{}{}{}}
\makeatother
\crefname{originallemma}{original lemma}{original lemmas}
\Crefname{originallemma}{Original lemma}{Original lemmas}
\makeatletter
\newlabel{lemma-quasi-compact-permanence}{{03KS}{source}{}{sourcepointer-3}{}}
\newlabel{lemma-quasi-compact-permanence@cref}{{[originallemma][0][]03KS}{[1][0][]source}{}{}{}}
\makeatother
\crefname{originallemma}{original lemma}{original lemmas}
\Crefname{originallemma}{Original lemma}{Original lemmas}
\makeatletter
\newlabel{lemma-reach-points-scheme-theoretic-image}{{089B}{source}{}{sourcepointer-4}{}}
\newlabel{lemma-reach-points-scheme-theoretic-image@cref}{{[originallemma][0][]089B}{[1][0][]source}{}{}{}}
\makeatother
\title{The valuative criterion characterizes universally closed quasi-compact quasi-separated morphisms of algebraic spaces}
\author{The Stacks Project Authors}
\date{stacks-project:spaces-morphisms.tex:local-snapshot-20261012; GNU FDL 1.2 or later}
\begin{document}
\maketitle
\setcounter{subsection}{1}
\noindent\textit{Editorial note.} All delivered mathematical text is the author's own, copied from the pinned Stacks Project chapter file spaces-morphisms.tex; no statement or proof step was rewritten, reordered or removed.\par\medskip
\noindent\textit{Editorial note.} The author states that a formal argument gives the reduction to the case of X\_A over Spec(A); the reduction is displayed as diagrams in the counted span and the paragraph that follows carries out the remaining work in full.\par\medskip
\noindent\textit{Editorial note.} The four references of the counted spans are recorded as non-load-bearing navigation pointers; the argument that consumes them stays inside the counted proof span.\par\medskip
\noindent\textit{Editorial note.} Preamble: the three native excerpt spans of sources/stacks-spaces-morphisms/preamble.tex (lines 8-13, 14-18 and 154-217). The chapter file itself carries no preamble text.\par\medskip
\noindent\textit{Allowed standard prerequisites (editorial).}\begin{itemize}
\item The converse criterion: quasi-compact plus valuative criterion implies universally closed. Lemma lemma-quasi-compact-existence-universally-closed, used at proof lines 8875-8877. Turns assertion (1) into assertion (2).
\item Reaching points of the scheme theoretic image of a morphism from a valuation ring. Lemma lemma-reach-points-scheme-theoretic-image, used at proof lines 8914-8922. Produces the extension A' of A dominating A with the closed point mapping to x'.
\item Permanence of quasi-compactness for the morphism x : Spec(K) -> X. Lemma lemma-quasi-compact-permanence, used at proof lines 8906-8908. Makes x quasi-compact, so that the closure of the image of x is the image of a closed set.
\item Base change stability of quasi-separatedness and universal closedness. used at proof lines 8894-8895. Justifies the formal reduction to the case Y = Spec(A).
\item The shape of the valuative criterion diagram. Definition definition-valuative-criterion, used at proof lines 8878-8893. Fixes the diagrams that have to be completed.
\end{itemize}\medskip
\ifdefined\hypertarget\hypertarget{sourcepointer-1}{}\fi
\noindent\textit{Source reference.} Original definition 03IX: focused-seeds/AGP/sources/stacks-snapshot-20261011/tags.txt, source lines 4619--4619. Navigation pointer into the Stacks Project; the standard fact it supplies is applied inside the counted proof span and the argument resting on it stays there.\par\medskip
\ifdefined\hypertarget\hypertarget{sourcepointer-2}{}\fi
\noindent\textit{Source reference.} Original lemma 03KA: focused-seeds/AGP/sources/stacks-snapshot-20261011/tags.txt, source lines 4670--4670. Navigation pointer into the Stacks Project; the standard fact it supplies is applied inside the counted proof span and the argument resting on it stays there.\par\medskip
\ifdefined\hypertarget\hypertarget{sourcepointer-3}{}\fi
\noindent\textit{Source reference.} Original lemma 03KS: focused-seeds/AGP/sources/stacks-snapshot-20261011/tags.txt, source lines 4688--4688. Navigation pointer into the Stacks Project; the standard fact it supplies is applied inside the counted proof span and the argument resting on it stays there.\par\medskip
\ifdefined\hypertarget\hypertarget{sourcepointer-4}{}\fi
\noindent\textit{Source reference.} Original lemma 089B: focused-seeds/AGP/sources/stacks-snapshot-20261011/tags.txt, source lines 10491--10491. Navigation pointer into the Stacks Project; the standard fact it supplies is applied inside the counted proof span and the argument resting on it stays there.\par\medskip
\begin{lemma}
\label{lemma-characterize-universally-closed-quasi-separated}
Let $S$ be a scheme. Let $f : X \to Y$ be a
morphism of algebraic spaces over $S$.
\begin{enumerate}
\item If $f$ is quasi-separated and universally closed, then
$f$ satisfies the existence part of the valuative criterion.
\item If $f$ is quasi-compact and quasi-separated, then
$f$ is universally closed if and only if the existence part of the
valuative criterion holds.
\end{enumerate}
\end{lemma}

\begin{proof}
If (1) is true then combined with
Lemma \ref{lemma-quasi-compact-existence-universally-closed}
we obtain (2). Assume $f$ is quasi-separated and universally closed.
Assume given a diagram
$$
\xymatrix{
\Spec(K) \ar[r] \ar[d] & X \ar[d] \\
\Spec(A) \ar[r] & Y
}
$$
as in Definition \ref{definition-valuative-criterion}.
A formal argument shows that the existence of the desired diagram
$$
\xymatrix{
\Spec(K') \ar[r] \ar[d] & \Spec(K) \ar[r] & X \ar[d] \\
\Spec(A') \ar[r] \ar[rru] & \Spec(A) \ar[r] & Y
}
$$
follows from existence in the case of the morphism $X_A \to \Spec(A)$.
Since being quasi-separated and universally closed are preserved
by base change, the lemma follows from the result in the next paragraph.

\medskip\noindent
Consider a solid diagram
$$
\xymatrix{
\Spec(K) \ar[r]_-x \ar[d] & X \ar[d]^f \\
\Spec(A) \ar@{=}[r] \ar@{..>}[ru] & \Spec(A)
}
$$
where $A$ is a valuation ring with field of fractions $K$.
By Lemma \ref{lemma-quasi-compact-permanence} and the fact
that $f$ is quasi-separated we have that
the morphism $x$ is quasi-compact.
Since $f$ is universally closed, we have in particular
that $|f|(\overline{\{x\}})$ is closed in $\Spec(A)$.
Since this image contains the generic point of $\Spec(A)$
there exists a point $x' \in |X|$ in the closure
of $x$ mapping to the closed point of $\Spec(A)$.
By Lemma \ref{lemma-reach-points-scheme-theoretic-image}
we can find a commutative diagram
$$
\xymatrix{
\Spec(K') \ar[r] \ar[d] & \Spec(K) \ar[d] \\
\Spec(A') \ar[r] & X
}
$$
such that the closed point of $\Spec(A')$ maps to $x' \in |X|$.
It follows that $\Spec(A') \to \Spec(A)$ maps the closed point
to the closed point, i.e., $A'$ dominates $A$ and this finishes the proof.
\end{proof}

\end{document}
